% Part: computability
% Chapter: computability-theory
% Section: k-1

\documentclass[../../../include/open-logic-section]{subfiles}

\begin{document}

\olfileid{cmp}{thy}{k1}
\olsection{د راکمېدنې يوه بېلګه}

راځئ د \olref[ppr]{prop:reduce} يوه کارونه وګورو۔

\begin{prop}
\ollabel{prop:k1}
فرض کړئ
\[
K_1 = \Setabs{e}{\cfind{e}(0) \fdefined}.
\]
بيا $K_1$ په محاسبوي ډول د شمېر وړ دے، خو محاسبه کېدونکے نۀ دے۔
\end{prop}

\begin{proof}
څرنګه چې $K_1 = \Setabs{e}{\lexists[s][T(e,0,s)]}$، نو د
\olref[eqc]{thm:exists-char} له مخې~$K_1$ په محاسبوي ډول د شمېر وړ
دے۔

د دې د ښودلو دپاره چې~$K_1$ محاسبه کېدونکے نۀ دے، راځئ وښيو چې
$K_0$ ورته راکمېدونکے دے۔

\begin{explain}
دا لږ ستونزمن دے، ځکه د~$K_1$ په کارولو سره يوازې د هغو محاسبو په
اړه پوښتنې کولے شو چې له ځانګړي ننوت، $0$، څخه پيلېږي۔ فرض کړئ يو
هوښيار ملګرے لرئ چې د دې ډول پوښتنو ځوابونه درکولے شي (داسې ملګري د
``اوراکلونو'' په نوم پېژندل کېږي)۔ بيا فرض کړئ څوک درته راشي او
وپوښتي چې~$\tuple{e, x}$ په~$K_0$ کښې دے که نه، يعنې ماشين~$e$ پر
ننوت~$x$ درېږي که نه۔ يو کار دا کولے شئ چې بل ماشين~$e_x$ جوړ کړئ؛
دا ماشين د \emph{هر} ننوت دپاره هغه ننوت له پامه غورځوي او پر ځاے ئې
~$e$ پر ننوت~$x$ چلوي۔ نو څرګنده ده چې د ماشين~$e$ پر ننوت~$x$
درېدل له دې پوښتنې سره معادل دي چې ماشين~$e_x$ پر ننوت~$0$ (يا هر
بل ننوت) درېږي که نه۔ بيا له خپل ملګري وپوښتئ چې دا نوی ماشين~$e_x$
پر ننوت~$0$ درېږي که نه؛ بدلې شوې پوښتنې ته د ملګري ځواب د اصلي
پوښتنې ځواب درکوي۔ دا له~$K_0$ څخه~$K_1$ ته موخه شوې راکمونه راکوي۔
\end{explain}

د نړيوالې جزوي محاسبه کېدونکې تابع په کارولو سره، $f$ هغه
درې‌ځاييزه تابع وي چې داسې تعريفېږي:
\[
f(x,y,z) \simeq \cfind{x}(y).
\]
پام وکړئ چې~$f$ خپل درېيم ننوت بيخي له پامه غورځوي۔ داسې شاخص~$e$
وټاکئ چې $f = \cfind{e}[3]$ وي؛ نو لرو:
\[
\cfind{e}[3](x,y,z) \simeq \cfind{x}(y).
\]
د $s$-$m$-$n$ قضيې له مخې داسې تابع~$s(e,x,y)$ شته چې د هر~$z$
دپاره
\begin{align*}
\cfind{s(e,x,y)}(z) & \simeq \cfind{e}[3](x,y,z) \\
& \simeq \cfind{x}(y).
\end{align*}

\begin{explain}
د پورته غير رسمي استدلال په ژبه، $s(e,x,y)$ د هغه ماشين شاخص دے چې
د هر ننوت~$z$ دپاره هغه ننوت له پامه غورځوي او~$\cfind{x}(y)$
محاسبه کوي۔
\end{explain}

په ځانګړي ډول، لرو:
\[
\cfind{s(e,x,y)}(0) \fdefined \quad \text{هله او يوازې هله چې} \quad
\cfind{x}(y) \fdefined.
\]
په بله وينا، $\tuple{x, y} \in K_0$ هله او يوازې هله چې
$s(e,x,y) \in K_1$ وي۔ نو هغه تابع~$g$ چې داسې تعريفېږي:
\[
g(w) = s(e,(w)_0,(w)_1)
\]
له~$K_0$ څخه~$K_1$ ته راکمونه ده۔
\end{proof}

\end{document}
